\section{Conclusion}
\label{sec:conclusion}

We began with a puzzle: can you identify a language model's architecture family from its
adversarial failures alone? Our answer is yes --- in principle, and at a large effect size ---
but statistical confirmation requires a larger model pool than the nine we evaluate here.

We introduced the \textbf{$\Delta^*$-vector framework} for multivariate architecture-family
adversarial profiling and provided the first formal effect-size measurement of between-family
$\Delta^*$ clustering: $\mathbf{\eta^2=0.293}$ (Cohen's $f^2 \approx 0.41$), exceeding the
pre-specified threshold in 83\% of attack categories. The strongest differentiation occurs on NLI
paraphrase attacks (adv\_rte, $\eta^2=0.592$), consistent with the hypothesis that bidirectional
and causal attention mechanisms process adversarial syntactic transformations differently. $N=9$
yields $\sim$40\% power; $N \geq 15$ is required for 80\% power and valid LOMO evaluation ---
a concrete design specification we provide as a contribution. A validated seven-module, 1,788-line
pipeline enables the community to replicate and extend this analysis.

\textbf{Future directions} include h-e1-v2 ($N=15$ for confirmatory study), alternative
classifiers (LDA, centroid) on existing $N=9$ data, attention concentration analysis
(h-m1--h-m4), CheckList behavioral partition, and large-scale (7B+) extension. The answer to
our opening puzzle is yes; the definitive proof awaits one well-powered follow-up study.
